
% Intriductory paragraph,
\subsection{Recursos Humanos}
Esta categoría comprende a las personas necesarias para ejecutar, guiar y validar la investigación:
\begin{itemize}
    \item \textbf{Investigador y tutor:} El investigador principal encargado del análisis, diseño, desarrollo y evaluación del sistema, trabajando en conjunto con la guía metodológica del tutor académico.
    \item \textbf{Personal clínico y administrativo:} Participación directa de un odontólogo general, un especialista y la secretaria del centro odontológico. El tiempo que estos profesionales destinan a las entrevistas, las observaciones de campo y la etapa de capacitación no representa un gasto monetario directo, sino un aporte valorizado esencial para la correcta adaptación del sistema a la realidad operativa.
\end{itemize}

\subsection{Recursos Institucionales}
Comprenden las entidades que brindan el marco académico y el espacio físico para el estudio:
\begin{itemize}
    \item \textbf{Institución de educación superior:} Provee las directrices formales de investigación, el acceso a repositorios bibliográficos y el respaldo para el desarrollo del proyecto.
    \item \textbf{Centro odontológico Bellidental:} Facilita sus instalaciones físicas y autoriza el acceso a su entorno de trabajo diario para efectuar el levantamiento de información y las mediciones de los ciclos de atención.
\end{itemize}

\subsection{Recursos Tecnológicos}
Agrupan los equipos y servicios digitales necesarios para construir e implementar la solución:
\begin{itemize}
    \item \textbf{Equipamiento informático:} Uso del computador personal del investigador, cuyo costo de desgaste prorrateado durante los meses de ejecución del proyecto se calcula en 200 dólares.
    \item \textbf{Herramientas de asistencia al desarrollo:} Suscripción a servicios de asistencia inteligente en la nube para el soporte continuo en la programación, por un valor de 20 dólares mensuales durante cuatro meses, totalizando 80 dólares.
    \item \textbf{Infraestructura en línea:} Alquiler de un servidor virtual básico por 30 dólares, el cual es necesario para alojar de forma centralizada y poner el sistema en funcionamiento durante la fase de pruebas e implementación final.
\end{itemize}

\subsection{Recursos Materiales y Operativos}
Cubren los elementos logísticos y físicos indispensables para la recolección de datos y el traslado:
\begin{itemize}
    \item \textbf{Instrumentos de análisis:} Acceso a los formatos actuales de fichas clínicas, hojas de abonos y recetarios físicos, cuyo uso es provisto sin costo por el establecimiento.
    \item \textbf{Movilización:} Visitas presenciales quincenales para el seguimiento y evaluación en la ciudad de operación de la clínica. Se calcula un total de ocho viajes de ida y vuelta con un costo de 3 dólares cada uno, sumando 24 dólares en total.
    \item \textbf{Suministros y fondo de imprevistos:} Un rubro de 16 dólares destinado a impresiones de guías de observación y gastos menores no planificados.
\end{itemize}

La suma de los rubros tecnológicos y logísticos detalla un presupuesto total requerido que asciende a 350 dólares. Esta cifra confirma la total viabilidad económica de la intervención, demostrando que se puede lograr una mejora sustancial en la reducción de los tiempos de gestión administrativa con un gasto económico prudente, realista y debidamente justificado.